\documentclass[12pt,a4paper]{article}
\usepackage{goyabase}

\title{\textbf{Laboratorio Avanzado de Patrones NoSQL: Cassandra}}
\author{Departamento de Informática}
\date{\today}

\usepackage[spanish, es-noshorthands, es-tabla]{babel}
\usepackage[autostyle]{csquotes}
\begin{document}

\maketitle

\section{Introducción al Laboratorio}\label{intro}
En este laboratorio exploraremos los patrones de diseño que hacen que gigantes como Netflix o Instagram usen Cassandra. No se trata solo de guardar datos, sino de modelarlos para que la lectura sea instantánea independientemente del volumen.

\begin{center}\rule{0.5\linewidth}{0.5pt}\end{center}

\subsection{Patrón 1: Series Temporales (Telemetría IoT)}\label{iot}
\textbf{Escenario:} Queremos guardar la temperatura de los sensores de los hoteles de GloboTour cada minuto.

\begin{tcolorbox}[title=Diseño de la Tabla]
\begin{lstlisting}[language=CQL]
CREATE TABLE sensor_readings (
    hotel_id uuid,
    sensor_id text,
    day date,
    reading_time timestamp,
    value float,
    PRIMARY KEY ((hotel_id, day), reading_time)
) WITH CLUSTERING ORDER BY (reading_time DESC);
\end{lstlisting}
\end{tcolorbox}

\textbf{Análisis del Patrón:}
\begin{itemize}
    \item \textbf{Composite Partition Key \texttt{(hotel\_id, day)}}: Agrupamos las lecturas por hotel \textbf{y por día}. Esto evita que una partición crezca demasiado con el paso de los años, manteniendo el rendimiento constante.
    \item \textbf{Uso Práctico:} Permite responder a \enquote{Dame todas las temperaturas de hoy del Hotel X} de forma ultra-eficiente.
\end{itemize}

\begin{center}\rule{0.5\linewidth}{0.5pt}\end{center}

\subsection{Patrón 2: Contadores Atómicos (Votos y Likes)}\label{counters}
\textbf{Escenario:} Necesitamos contar cuántas veces se ha visualizado cada oferta de viaje sin que las escrituras masivas bloqueen la tabla.

\begin{tcolorbox}[title=Uso de la tabla \texttt{offer\_stats}]
\begin{lstlisting}[language=CQL]
-- Nota: Las tablas de contadores solo pueden tener contadores y la PK.
UPDATE offer_stats 
SET views = views + 1 
WHERE offer_id = 'OFERTA-PARIS-2024';
\end{lstlisting}
\end{tcolorbox}

\textbf{Análisis del Patrón:}
Cassandra gestiona los contadores de forma especial para evitar \enquote{Race Conditions} sin necesidad de bloqueos de fila de SQL.

\begin{center}\rule{0.5\linewidth}{0.5pt}\end{center}

\subsection{Patrón 3: Datos con Caducidad (TTL)}\label{ttl}
\textbf{Escenario:} Marketing lanza ofertas \enquote{Flash} que solo duran 2 horas. No queremos un proceso que borre datos viejos (que es lento); queremos que los datos \textbf{mueran solos}.

\begin{tcolorbox}[title=Inserción con TTL]
\begin{lstlisting}[language=CQL]
INSERT INTO flash_offers (offer_id, discount) 
VALUES ('REBAJA-10', 0.15) 
USING TTL 7200; -- Expira en 2 horas (7200 seg)
\end{lstlisting}
\end{tcolorbox}

\begin{center}\rule{0.5\linewidth}{0.5pt}\end{center}

\subsection{Laboratorio de Depuración: ¿Por qué falla mi CQL?}\label{debug}

\begin{itemize}
    \item \textbf{Error de Filtrado:} \texttt{SELECT * FROM sensor\_readings WHERE value > 25;} $\rightarrow$ \textbf{FALLARÁ}. Cassandra no sabe en qué nodo está ese valor sin la clave de partición.
    \item \textbf{Error de Ordenación:} \texttt{SELECT * FROM sensor\_readings ORDER BY value DESC;} $\rightarrow$ \textbf{FALLARÁ}. Solo puedes ordenar por las \textit{Clustering Columns} definidas en el esquema.
\end{itemize}

\begin{tcolorbox}[title=La Regla de Oro, colframe=orange!75!black]
En Cassandra, si no puedes responder a una pregunta con tu tabla actual, \textbf{crea otra tabla}. El disco es barato, el tiempo de tus usuarios no.
\end{tcolorbox}

\end{document}
